\section{A quantitative discount reduction}
\label{sec:discount}

We first reduce the sign of each mean-payoff value to an approximation of
one discounted fixed point. The reduction needs two quantitative facts:
nonzero values are separated from zero, and a discount of polynomial bit
length approximates every value more accurately than this separation.
We prove both facts directly, including the identification of the limiting
discounted value with the expected pathwise lower average.

The discounted fixed-point method originates in Shapley's work
\cite{Shapley1953}. Our argument for a single positional pair at all small
discounts is a quantitative version of the eventual-optimality principle
introduced by Blackwell for finite Markov decision processes
\cite{Blackwell1962}; stationary optimality for perfect-information
stochastic games with limiting-average payoffs was established by Liggett
and Lippman \cite{LiggettLippman1969}. Quantitative thresholds for
small-discount optimality are studied further in
\cite{GaubertGrandClementKatz2025}. The estimates and the payoff
convention needed here are proved below.

Let $n=|V|$, and define
\begin{equation}
\begin{split}
 W&=\max\bigl(1,\max_{e\in E}|w(e)|\bigr),\\
 Q&=\prod_{\substack{e\text{ outgoing from}\\i\in V_{\mathrm{ch}}}}
       \operatorname{den}(p_i(e)),\\
 A_0&=n!\,[Q(3+W)]^n,
 \qquad H=(W+4)A_0,
 \qquad \lambda_0=\frac1{32H^2}.
\end{split}
\label{eq:discount-parameters}
\end{equation}
Here $W,Q,A_0,H$ are positive integers. The denominator in the product
defining $Q$ is the positive denominator supplied in the input; the
fractions need not be reduced. An empty product is one.
The explicit input gives $n=O(L)$ and $\log Q,\log W=O(L)$, because
the denominator and reward bits are included in $L$. Therefore
\[
 \log H=O\!\left(n\bigl(\log(n+1)+\log Q+\log(W+3)\bigr)\right)
        =O((L+2)^2),
\]
so $\lambda_0$ has polynomial bit length.

For $0<u<1$, define $f_u:\R^V\to\R^V$ by taking at vertex $i$ the
maximum, minimum, or prescribed chance average of
\begin{equation}
                 u w(e)+(1-u)t_j,
                 \qquad e:i\longrightarrow j,
\label{eq:discount-edge}
\end{equation}
according as $i$ belongs to Max, Min, or Chance. This map is monotone and
is a contraction in the maximum norm, with constant $1-u$. Indeed, each
expression in \eqref{eq:discount-edge} changes by at most
$(1-u)\|t-t'\|_\infty$, and taking an extremum or an average preserves
that bound. Its iterates from zero converge, because the successive
differences decrease geometrically. Their limit is the unique fixed
point, denoted $v(u)$. Since $f_u$ preserves $[-W,W]^V$,
\begin{equation}
                       \|v(u)\|_\infty\le W.
\label{eq:discount-bound}
\end{equation}

\begin{theorem}[Quantitative discount reduction]
\label{thm:discount}
For every game in the input class, there are vectors $G,\psi\in\R^V$
and a pure positional strategy pair $(\sigma^*,\tau^*)$ whose selected
edges attain the extrema defining $f_u(v(u))$ for every
$0<u\le1/(2H)$. Moreover,
\begin{equation}
                 v_i(u)=G_i+u\psi_i+O(u^2)
                 \qquad(u\downarrow0).
\label{eq:gain-bias-expansion}
\end{equation}
For every start vertex $i$ and every behavioral opposing strategy,
\begin{equation}
 \E_{i,\sigma^*,\tau}[\operatorname{MP}_w]\ge G_i,
 \qquad
 \E_{i,\sigma,\tau^*}[\operatorname{MP}_w]\le G_i.
\label{eq:gain-strategy-bounds}
\end{equation}
Consequently $G_i=\val(i)$. Each $G_i$ is rational with a representation
whose denominator has absolute value at most $A_0$, and
\begin{equation}
 \begin{gathered}
 |G_i|\le W,
 \qquad G_i\ne0\ \Longrightarrow\ |G_i|\ge\frac1H,\\
 |v_i(\lambda_0)-G_i|\le\frac1{16H}.
 \end{gathered}
\label{eq:discount-gap-error}
\end{equation}
\end{theorem}

The proof has an algebraic part and a probabilistic part. Rational
functions first give a pair that works for every sufficiently small
discount, together with its limiting gain $G$ and first-order term
$\psi$. We then use these two vectors to control individual plays
against arbitrary behavioral opponents.

\subsection{A positional pair valid for all small discounts}

Temporarily fix a pure positional pair: one outgoing edge at each Max
and Min vertex. Let $P$ be its row-stochastic transition matrix and let
$r_i$ be its expected one-step reward at $i$. At a chance vertex the
transition probabilities of parallel edges are added in $P$, whereas
their possibly different rewards are averaged in $r$. The pair's
discounted equations are
\begin{equation}
                       x(u)=u r+(1-u)P x(u).
\label{eq:pair-discount-equations}
\end{equation}
The matrix $I-(1-u)P$ is invertible for $0<u<1$: a vector in its kernel
would satisfy $\|x\|_\infty\le(1-u)\|x\|_\infty$. The corresponding
affine contraction preserves $[-W,W]^V$, so its solution satisfies
$\|x(u)\|_\infty\le W$.

The determinant argument follows the rational-function sign method
of Andersson and Miltersen \cite[Lemma~1]{AnderssonMiltersen2009};
we derive all coefficient bounds in the input parameters above.
For a polynomial $a$, write $\|a\|_{\mathrm{coef}}$ for the sum of
the absolute values of its coefficients. This norm is subadditive and
submultiplicative. Every entry of $QP$ and $Qr$ is an integer. Cramer's
rule therefore gives integer polynomials $N_i,D$ with
\begin{equation}
 x_i(u)=\frac{N_i(u)}{D(u)},
 \qquad D(u)=\det\bigl(Q(I-(1-u)P)\bigr).
\label{eq:discount-cramer}
\end{equation}
Each entry of the matrix defining $D$ has coefficient norm at most $3Q$.
The replacing column defining $N_i$ is $Qu r$, whose entries have
coefficient norm at most $QW$. Expanding each determinant over its $n!$
permutations gives
\begin{equation}
 \|D\|_{\mathrm{coef}}\le A_0,
 \qquad \|N_i\|_{\mathrm{coef}}\le A_0.
\label{eq:discount-polynomial-bound}
\end{equation}

Whether this pair's chosen edge is optimal at a player vertex $i$ is
determined by comparisons with every outgoing edge $e:i\to j$. Each
comparison is the sign of
\begin{equation}
 x_i(u)-u w(e)-(1-u)x_j(u)
   =\frac{N_i(u)-u w(e)D(u)-(1-u)N_j(u)}{D(u)}.
\label{eq:discount-edge-comparison}
\end{equation}
The numerator has coefficient norm at most $(3+W)A_0\le H$.

To control these signs, let $a$ be any nonzero integer polynomial with
$\|a\|_{\mathrm{coef}}\le H$. Factoring its lowest power of $u$ gives
$a(u)=u^\ell(c+u q(u))$, where $c$ is a nonzero integer. For
$0<u\le1/(2H)$,
\[
                   |u q(u)|\le Hu\le\tfrac12<|c|.
\]
Thus the sign of $a(u)$ throughout this interval is the sign of its
lowest nonzero coefficient. This observation applies to both the
numerator and denominator in \eqref{eq:discount-edge-comparison};
the denominator is nonzero there by invertibility. Every edge comparison
therefore has a constant sign or vanishes identically on the interval.

Now choose attaining edges in the maximum and minimum defining
$f_{\lambda_0}(v(\lambda_0))$. They give a pure positional pair
$(\sigma^*,\tau^*)$, and its linear solution at $\lambda_0$ equals
$v(\lambda_0)$. Since $\lambda_0\le1/(2H)$, all edge inequalities
that hold at $\lambda_0$ continue to hold for every $0<u\le1/(2H)$.
For each such $u$, the pair's linear solution is consequently a fixed
point of $f_u$, hence equals $v(u)$. The pair is needed for the proof;
the algorithm need not compute it.

\subsection{The rational gain and its approximation error}

Use the polynomials of the pair just selected. Let $\ell$ be the lowest
nonzero degree of $D$. Boundedness of $N_i(u)/D(u)$ as $u\downarrow0$
implies that $u^\ell$ also divides each $N_i$. Define
$\widehat D=D/u^\ell$ and $\widehat N_i=N_i/u^\ell$.
Their coefficient norms satisfy \eqref{eq:discount-polynomial-bound},
and $\widehat D(0)$ is a nonzero integer. The quotient is analytic at
zero, so it has the expansion \eqref{eq:gain-bias-expansion}, with
\begin{equation}
                         G_i=\frac{\widehat N_i(0)}{\widehat D(0)}.
\label{eq:rational-gain}
\end{equation}
The bound \eqref{eq:discount-bound} gives $|G_i|\le W$. The displayed
denominator has magnitude at most $A_0$, and its numerator is an
integer. Thus $G_i\ne0$ implies $|G_i|\ge1/A_0\ge1/H$.

The same coefficient bounds also control the approximation error.
The polynomial $\widehat N_i-G_i\widehat D$ has constant coefficient
zero and coefficient norm at most $(1+W)A_0$. For
$0<u\le1/(2H)$, the tail of $\widehat D$ has magnitude at most
$A_0u\le1/2$, so $|\widehat D(u)|\ge1/2$. It follows that
\begin{equation}
 |v_i(u)-G_i|
   \le 2(1+W)A_0u
   \le 2Hu.
\label{eq:uniform-discount-error}
\end{equation}
At $u=\lambda_0$, this is the error claimed in
\eqref{eq:discount-gap-error}.

It remains to identify $G_i$ with the mean-payoff value. The gain
controls which transitions can persist, and the coefficient $\psi$
then controls the rewards on those transitions. For a selected player
edge $e:i\to j$, substitution of \eqref{eq:gain-bias-expansion} into
the discounted equation gives
\begin{equation}
               G_i=G_j,
               \qquad \psi_i=w(e)+\psi_j-G_j.
\label{eq:selected-gain-bias}
\end{equation}
At a chance vertex, the corresponding identities are
\begin{equation}
 \begin{split}
 G_i&=\sum_{e:i\to j}p_i(e)G_j,\\
 \psi_i&=\sum_{e:i\to j}p_i(e)\bigl(w(e)+\psi_j-G_j\bigr).
 \end{split}
\label{eq:chance-gain-bias}
\end{equation}
At a Min vertex, the inequality
$v_i(u)\le u w(e)+(1-u)v_j(u)$ for an arbitrary outgoing edge first
gives $G_i\le G_j$. If $G_i=G_j$, subtracting this common constant,
dividing by $u$, and letting $u\downarrow0$ gives
\begin{equation}
                  \psi_i\le w(e)+\psi_j-G_j.
\label{eq:min-gain-bias}
\end{equation}
At a Max vertex, the gain inequality and, when the gains agree, the
bias inequality both reverse.

\subsection{From one-step inequalities to pathwise payoffs}

We use two elementary facts about adapted processes. Fixing a selected
strategy will give the gain along a play one-sided conditional drift.
The first lemma shows that the gain then becomes constant, so only
gain-preserving player moves can persist and the bias inequalities
apply. The second lemma shows that the centered contributions at chance
vertices have vanishing averages.

\begin{lemma}[Finite-valued submartingales]
\label{lem:finite-submartingale}
Let $(Z_t)_{t\ge0}$ be adapted to a filtration $(\mathcal F_t)_{t\ge0}$
and take values in a fixed finite subset of $\R$. If
$\E[Z_{t+1}\mid\mathcal F_t]\ge Z_t$ for every $t$, then $Z_t$ is
eventually constant almost surely. Its eventual value $Z_\infty$
satisfies $\E Z_\infty\ge\E Z_0$.
\end{lemma}

\begin{proof}
Add a constant to make $Y_t=Z_t+c$ nonnegative. Conditional expectation
and the assumed drift give
\[
 \begin{split}
 \E[Y_{t+1}^2-Y_t^2]
 &=\E[(Y_{t+1}-Y_t)^2]
   +2\E\!\left[Y_t\,\E[Y_{t+1}-Y_t\mid\mathcal F_t]\right]\\
 &\ge\E[(Y_{t+1}-Y_t)^2].
 \end{split}
\]
Summing and using boundedness shows that
$\E\sum_{t\ge0}(Y_{t+1}-Y_t)^2<\infty$. The sum is therefore finite
almost surely. If there is more than one possible value, let $\Delta>0$
be the minimum distance between distinct values. Every change contributes
at least $\Delta^2$ to the sum, so only finitely many changes occur.
The singleton case is immediate. Finally, $\E Z_t\ge\E Z_0$ for all
$t$, and bounded convergence gives the claimed expectation inequality.
\end{proof}

\begin{lemma}[Averages of bounded martingale differences]
\label{lem:bounded-martingale-average}
Let $(\mathcal F_t)_{t\ge0}$ be a filtration. Suppose $D_t$ is
$\mathcal F_{t+1}$-measurable,
$\E[D_t\mid\mathcal F_t]=0$, and $|D_t|\le c$ for a fixed finite
constant $c$. Then $T^{-1}\sum_{t<T}D_t\to0$ almost surely.
\end{lemma}

\begin{proof}
For $s<t$, conditioning on $\mathcal F_t$ gives $\E[D_sD_t]=0$.
Hence $\E[(\sum_{t<T}D_t)^2]\le Tc^2$. Chebyshev's inequality at
$T=j^2$ gives, for every $\varepsilon>0$,
\[
 \mathbb P\!\left(\left|\frac1{j^2}\sum_{t<j^2}D_t\right|>
          \varepsilon\right)
 \le\frac{c^2}{\varepsilon^2j^2}.
\]
The sum over $j$ is finite. The probability that one of these events
occurs after index $m$ is at most the tail of this sum, which tends to
zero. Applying this argument to $\varepsilon=1,1/2,1/3,\ldots$ proves
convergence along squares on a single event of probability one.
For $j^2\le T<(j+1)^2$, the additional partial sum has magnitude at
most $c(2j+1)$. Dividing by $T\ge j^2$ extends convergence to every
$T$.
\end{proof}

\begin{proof}[Completion of the proof of Theorem~\ref{thm:discount}]
Fix a start vertex $i$, fix Max to $\sigma^*$, and let Min use an
arbitrary behavioral strategy. Write $X_t$ for the vertex before the
edge $e_t$ is traversed, and let $\mathcal F_t$ be the sigma-field
generated by the complete observed history through $X_t$. At a chance
vertex the conditional law of $e_t$ is the prescribed distribution.
At a player vertex it is the
distribution induced by that player's strategy, including any private
randomness.

The process $G_{X_t}$ has nonnegative conditional drift. A selected Max
edge preserves $G$ by \eqref{eq:selected-gain-bias}; every Min edge
weakly increases $G$; and the first identity in
\eqref{eq:chance-gain-bias} gives zero drift at a chance vertex.
Lemma~\ref{lem:finite-submartingale} gives an eventual value
$G_\infty$ with
\begin{equation}
                         \E G_\infty\ge G_i.
\label{eq:gain-limit-expectation}
\end{equation}

Define the one-step residual and its chance part by
\[
 R_t=w(e_t)+\psi_{X_{t+1}}-\psi_{X_t}-G_{X_t},
 \qquad
 D_t=\mathbf 1_{\{X_t\in V_{\mathrm{ch}}\}}R_t.
\]
These random variables are uniformly bounded because the arena is
finite. At a chance vertex, the two identities in
\eqref{eq:chance-gain-bias} imply
\[
 \E[R_t\mid\mathcal F_t]
   =\sum_{e:X_t\to j}p_{X_t}(e)(w(e)+\psi_j)
      -\psi_{X_t}-G_{X_t}=0.
\]
Outside chance vertices $D_t=0$. Thus
$\E[D_t\mid\mathcal F_t]=0$, and
Lemma~\ref{lem:bounded-martingale-average} gives
$T^{-1}\sum_{t<T}D_t\to0$ almost surely.

After $G_{X_t}$ has become constant, every traversed nonchance edge
preserves $G$. Its residual is zero at Max vertices by
\eqref{eq:selected-gain-bias}, and nonnegative at Min vertices by
\eqref{eq:min-gain-bias}. The nonchance residuals are therefore
nonnegative except for a finite initial segment. Consequently
\[
                    \liminf_{T\to\infty}\frac1T\sum_{t<T}R_t\ge0
                    \qquad\text{almost surely}.
\]
Telescoping the bias terms gives
\begin{equation}
 \frac1T\sum_{t<T}w(e_t)
   =\frac1T\sum_{t<T}G_{X_t}
     +\frac1T\sum_{t<T}R_t
     +\frac{\psi_{X_0}-\psi_{X_T}}{T}.
\label{eq:gain-bias-telescoping}
\end{equation}
The first average tends to $G_\infty$, and the last term tends to zero.
Hence $\operatorname{MP}_w\ge G_\infty$ almost surely. Taking
expectations and using \eqref{eq:gain-limit-expectation} proves the
first inequality in \eqref{eq:gain-strategy-bounds}.

For the second inequality, fix Min to $\tau^*$ and let Max be arbitrary.
Now $-G_{X_t}$ has nonnegative conditional drift. The same finite-valued
lemma yields an eventual gain $G_\infty$ with $\E G_\infty\le G_i$.
The chance residuals still have average zero. At Min vertices the
selected edges have residual zero; at Max vertices an edge preserving
$G$ has nonpositive residual by the reversed version of
\eqref{eq:min-gain-bias}. Thus the nonchance residuals are eventually
nonpositive, and \eqref{eq:gain-bias-telescoping} gives
\[
             \limsup_{T\to\infty}\frac1T\sum_{t<T}w(e_t)
                    \le G_\infty
                    \qquad\text{almost surely}.
\]
This also bounds the pathwise lower average. Taking expectations proves
the second inequality in \eqref{eq:gain-strategy-bounds}.

Max can therefore guarantee $G_i$, and Min can guarantee that the payoff
is at most $G_i$, so $\val(i)=G_i$. The payoffs and gain limits in these
expectation comparisons are bounded. The argument permits each opposing
history-dependent distribution and establishes the expected-liminf
convention directly.
\end{proof}
